\section{The exceptional primes of the selected hub}

% BEGIN X78A1
\paragraph{X78A1. The integer windows lie strictly between zero and three times the scale.}
Fix a size datum $L$, its parameter choice $P$, and the original
$\mathrm{SizesAt}(L)$ inequalities. Write $U=x^\varepsilon$,
$b=e^U$, $X'=\exp(x^{\varepsilon/3})$, $b'=xb/X'$,
$w''=b'/(xn)$, $S_c=L.\mathrm{Sc}$, $\eta=L.\mathrm{etaSz}$ and
$p_*=\ell_*(L)$. The ranks satisfy $m\ge2$, $t>3m$ and $n\ge1$.
For $s=\lfloor 2b/p_*\rfloor$ the centers and integer windows are
\[
 \begin{aligned}
 \mathcal B={}&\{b,w''\}\cup
    \{b_\tau+\lfloor b'/x\rfloor:\tau\in\mathrm{TT}\},\\
 I_W={}&\mathbb Z\cap
 [W-\eta-2s-3,\ W+S_c+2\eta+2s+4].
 \end{aligned}
\]
Every $a\in I_W$ for $W\in\mathcal B$ satisfies $0<a<3b$,
and $p_*^2>8b$.

Indeed $x>1$, $U\ge1$ and the denominator of
\[
 p_*=\frac{mb}{nt^2(N(m)+t)U^2+kn^2(S_c+t+U)}
\]
is at least $t>2m$. Thus $0<2p_*<b$. The two defining size
inequalities at $p_*$ give
\[
 b'+S_c+\eta+b/p_*<p_*,
 \qquad 2\eta+t+8b/p_*+7<p_* .
\]
The second implies $p_*^2>8b$. Since $0\le s\le 2b/p_*$,
the first gives $S_c<p_*$ and the second gives
$2\eta+2s+4<p_*$. Hence the upper offset of every window is
strictly less than $2p_*<b$.

For its lower endpoint, $w''=b/(nX')\le b$ and $w''\le b'$.
The center indexed by $\tau$ is at least $b'$ by the input
$b_\tau\ge b'$ and $\lfloor b'/x\rfloor\ge0$. Thus every
center is at least $w''$. The parameter conditions give
\[
 w''\ge m(t+1)\eta\ge16\eta,\qquad
 12b/p_*+2t+8\le\eta/m\le\eta.
\]
In particular $2s+3<\eta$, proving positivity throughout the windows.

For their upper endpoints put $A=\mathrm{szNst}(L)U L.\mathrm{LaG}/k$.
The size bound $\mathrm{SizeRho}$ gives $4A\le x^{-2}$, since
its other term is nonnegative by $L.\mathrm{LaG}\ge U$ and $x>1$.
Then $A\le1/4$, and $\mathrm{SizeBot}$ gives
$b_\tau\le (1+A)b'\le5b'/4$. Since $\lfloor b'/x\rfloor\le b'$
and $b'<p_*<b/2$, each such center is at most $9b'/4<2b$.
The other two centers are at most $b$. Adding the upper offset
proves the claimed strict upper bound.
% END X78A1

% BEGIN X78A2
\paragraph{X78A2. The exact excluded set and both of its bounds.}
In the setting of X78A1, form $E_L$ by taking, for each
$W\in\mathcal B$ and $a\in I_W$, all prime divisors of $a$
strictly greater than $p_*$, and taking the union. This is a finite
set of natural numbers. By positivity of $a$, divisibility in
$\mathbb Z$ by a positive rational prime is equivalent to divisibility
of its natural absolute value. Thus its underlying set is exactly
$\mathrm{BaseSizes.Exc}(L)$, including the strict threshold.

At most one such prime divides any given $a$. Indeed distinct
prime divisors $p,q>p_*$ would be coprime with $pq\mid a$, whereas
$0<a<3b<8b<p_*^2<pq$. Every prime in $E_L$ is itself at most
$a<3b$ for one of these integers.

The number of integers between real endpoints $u\le v$, inclusive,
is at most $v-u+1$, by taking the ceiling and floor. Each window
therefore has at most $S_c+3\eta+4s+8$ elements. There are at most
$|\mathrm{TT}|+2$ distinct centers, including when some coincide.
Summing the bounds for the finite union and using $4s\le8b/p_*$
gives
\[
 |E_L|\le (|\mathrm{TT}|+2)
       (S_c+3\eta+t+8b/p_*+9),\qquad
 p\in E_L\ \Longrightarrow\ p\le3b .
\]
These are both clauses of the original $\mathrm{ExcOK}(L,E_L)$.
% END X78A2

% BEGIN X78A3
\paragraph{X78A3. The same exceptional set in the family, setup and table.}
For every original family $\mathrm{Fm}$ and $x>\mathrm{Fm}.x_1$,
each $i\in\mathrm{Fm}.\mathrm{Ix}(x)$ supplies $L=\mathrm{Fm}.L(x,i)$,
its parameter choice and $\mathrm{SizesAt}(L)$. Apply X78A2 to
these data. The original uniform assertion
$\mathrm{excludedPrimes\_few\_and\_small\_along\_family}$
therefore admits the explicit cutoff $x_0=\mathrm{Fm}.x_1$.

Under the hypotheses of F78D2, keep its cutoff. At each permitted
$x,i$ with the prescribed odd Tate--Shafarevich finiteness, choose
$E_L$ first and apply F78D2 to that finite set. It returns the
compatible Wach records and a hub $X$ with
\[
 X.\mathrm{exc}=E_L,\qquad
 (X.\mathrm{exc}:\mathrm{Set}\ \mathbb N)=\mathrm{BaseSizes.Exc}(L),
 \qquad \mathrm{ExcOK}(L,X.\mathrm{exc}).
\]
Hence the actual $\mathrm{Concl.exc}$ field holds for the same $X$
as the setup and killing predicate. When $C_2\ge1$ and
$C_s\ge\max(3C_2+1,22365)$, F78D3 supplies the five numerical
clauses on this $X$ as well. The equality with
$\mathrm{BaseSizes.Exc}(L)$ provides the exceptional-set identity
required by the large-prime congruence argument.
% END X78A3
